\section{RQ3: Are the Changes Benign or Risky?}\label{sec:rq3}

Validation (Section~\ref{sec:validation}) establishes which flagged changes are real. This section asks which of the real ones matter, using the same 80 labeled events.

\paragraph{What security-relevant changes look like}
Of the 43 real changes, 13 (30\%) are security-relevant; the rest are routine for the tool's purpose, such as reading a non-secret configuration variable, writing the server's own cache, or adding a log file. The security-relevant ones fall into four groups. Six concern credentials: device login passwords and a storage service's secret key read from the environment, an internal token-minting secret, a map of per-account API keys, an identity-provider token exchange that uses the caller's bearer token, and environment-held credentials in a ticketing tool. Three let the AI decide what runs or where data is written: a test-case listing tool that gained an optional parameter naming the file to overwrite, a data-analysis server that started executing model-written Python inside the user's live session, and a tool whose output is a shell command its description tells the agent to run. Three add a new local channel or store: a loopback WebSocket that replaces the browser's native-messaging bridge, a database file created by a permissions tool in a server that lets agents build their own tools, and a voice-control switch in a desktop-automation server. The last is a description that now directs the agent to move saved credentials into a live web action, with explicit safeguards.

\begin{finding}{3}
Most real changes are routine, and that is the problem. About 30\% are security-relevant: new access to credentials, AI-controlled execution or file writes, and new local channels. Risky changes are a minority inside a steady flow of routine updates.
\end{finding}
